\section{Offenbarung}

Gott hat unendliches und vollkommens Wissen, während der Mensch nur endliches und unvollkommenes Wissen besitzt. Der Mensch kann unabhängig von der Schrift nicht einmal völlig wissen, was die Schöpfung offenbart.

In der Schrift gibt es zwei Offenbarungsformen:
\begin{enumerate}
    \item die allgemeine Offenbarung \biblezit{Ps}{19:2-7}
    \item die besondere Offenbarung \biblezit{Ps}{19:8-15}
\end{enumerate}

\begin{otsblock}{Allgemeine Offenbarung}
\begin{bibelbox}{SCHL}{Ps}{19:2-7}
    \hh{2}Die Himmel erzählen die Herrlichkeit Gottes, und die Ausdehnung verkündigt das Werk seiner Hände.

    \hh{3}Es fliesst die Rede Tag für Tag. Nacht für Nacht tut sich die Botschaft kund.

    \hh{4}Es ist keine Rede und es sind keine Worte, deren Simme unhörbar wäre.

    \hh{5}Ihre Reichweite erstreckt sich über die ganze Erde, und ihre Worte bis ans Ende des Erdkreises. Er hat der Sonne am Himmel ein Zelt gemacht.

    \hh{6}Und sie geht hervor wie ein Bräutigam aus seiner Kammer und freut sich wie ein Held, die Bahn zu durchlaufen.

    \hh{7} Sie geht an einem Ende des Himmels auf und läuft um bis ans andere Ende, und nichts bleibt vor ihrer Glut verborgen.
    \textit{Interessant \enquote{ihrer Glut}}
\end{bibelbox}
Die allgemeine Offenbarungist Gottes Zeugnis über sich selbst, und zwar mittels der Schöpfung an seine Geschöfe. Wenn jemand zum Himmel schaut, bezeugt das Universum selbest die Tatsache, dass es einen Schöpfer gibt.

\begin{bibelbox}{SCHL}{Ps}{19:2}
Die Himmel erzählen die Herrlichkeit Gottes, und die Ausdehnung verkündet seiner Hände Werk.
\end{bibelbox}
Je mehr jemand die unendlichen Weiten des Weltalls oder auch die kleinsten Partikel der Molekularstruktur erforscht, desto mehr ist er gedrängt, mit Staunen und Ehrfurcht die wahre Grösse der Schöpfers anzuerkennen.

\betonung[b, p]{Beispiele von Wissenschaftler}

Eine ander Form der allgeimeinen Offenbarung ist das Erkennen von Richtig und Falsch und das Werk des Gewissens, das Sünder anklagt.
\begin{bibelbox}{SCHL}{Rom}{2:14-15}
Wenn nähmlich Heiden, die das Gesetz nicht haben, doch von Natur aus tun, was das Gesetz verlangt, so sind sie, die das Gesetz nicht haben, sich selbst ein Gesetz,\\
da sie ja beweisen, dass das Werk des Gesetzes in ihren Herzen geschrieben ist, was auch ihr Gewissen bezeugt, dazu ihre Überlegungen, die sich untereinander verklagen oder auch entschuldigen --
\end{bibelbox}
Jeder Mensch beginnt mit dem innerlichen Begreifen der Tatsache, dass der Mensch zwar endlich ist, dass es aber in seiner Existens um mehr geht als nur diese zeitlich begrenzte Wirklichkeit.

Während die allgemeine Offenbarung sehr viel von der Macht, Weisheit, Güte, Gerechtigkeit und Majestät des Schöpfers vermittelt, ist auf das beschränkt, was durch den sündigen Menschen beobachtet werden kann. Das letztendliche Ziel der allgemeinden Offenbarung ist es, den Menschen keine Entschuldigung dafür zu lassen, dass sie das Wesen ihres Schöpfers nicht erkannt und anerkannt haben.

Niemand kann durch die allgemeine Offenbarung errettet werden. \biblezit{Rom}{10:5-17}, \biblezit{IKor}{1:18-2:5}

Das Problem des Menschen ist also nicht, dass ihm Informationen fehlen, sondern die fehlende Bereitschaft, diese Informationen anzuerkennen und danach zu handeln.
\end{otsblock}
